\documentclass[10pt,a4paper]{amsart}
\usepackage[margin=19mm]{geometry}
\usepackage{amsmath,amssymb,mathtools}
\usepackage[scr=rsfs,cal=euler]{mathalfa}
\usepackage{xfrac}
\usepackage[unicode,hidelinks,bookmarks=false]{hyperref}
\newcommand{\ie}{i.e.\ }
\newcommand{\cf}{cf.\ }
\newcommand{\resp}{respectively\ }
\newcommand{\Subsection}{\subsection}
\providecommand{\nobreakdash}{\nobreak\hbox{-}}
\setlength{\emergencystretch}{2em}
\multlinegap0pt
\usepackage{bm}
\usepackage{bbm}
\usepackage{dsfont}
\usepackage{xspace}
\usepackage{cancel}
\usepackage{mathtools}
\usepackage{amsthm}
\makeatletter
\@ifundefined{lemma}{\newtheorem{lemma}{Lemma}}{}
\@ifundefined{proposition}{\newtheorem{proposition}{Proposition}}{}
\makeatother
\newcommand{\KS}{\mathrm{KS}}
\newcommand{\LB}{\mathrm{LB}}
\newcommand{\bg}{\mathrm{bg}}
\newcommand{\refi}{\mathrm{ref}}
\title[A certified refinement and asymptotic analysis of the Kuznet]{A certified refinement and asymptotic analysis of the Kuznetsov-Sahinidis diameter bound for Lennard-Jones clusters}
\author{Guillaume Lecomte}
\date{Source: arXiv:2607.09555v1 (CC BY 4.0)}
\begin{document}
\maketitle

\noindent\emph{Source and licence.} A certified refinement and asymptotic analysis of the Kuznetsov-Sahinidis diameter bound for Lennard-Jones clusters. Version arXiv:2607.09555v1 (CC BY 4.0), distributed under the Creative Commons Attribution 4.0 International licence, as recorded in the arXiv licence metadata for this exact version and shown on the abstract page https://arxiv.org/abs/2607.09555v1. Full rights notice: licenses/arxiv-2607.09555-CC-BY-4.0.txt.\par\smallskip

\noindent\textit{Editorial scope.} Counted unit: the Proposition whose source label is \texttt{prop:gap} in the article (source0.tex lines 362--372, source label \texttt{prop:gap}), delivered together with the article's own complete original proof of it (source0.tex lines 374--404, one \texttt{proof} environment of 31 lines). No mathematics is written, repaired or re-derived: statement and proof are the article's own text line for line. Required context retained uncounted: eq:kernel (lines 71--77); eq:elb (lines 81--83); eq:f (lines 93--101); eq:elbsplit (lines 285--287); eq:square (lines 295--297); eq:minelb (lines 299--301); lem:penalty (lines 350--356). Every definition, display and label of those lines is retained unchanged. Omitted from this package and not redistributed: 1--70, 78--80, 84--92, 102--284, 288--294, 298--298, 302--349, 357--361, 373--373, 405--491. Nothing inside a retained excerpt is omitted or altered: every retained body is delivered exactly as the article's lines are written, so no omission receipt of any kind accompanies this package. Citations: the retained excerpts carry no citation command at all, so no bibliography entry of the article is transcribed and no reference is redistributed. Reference adaptation: none was needed and none was made; every reference inside the delivered unit resolves inside it, so no unresolved reference or citation is produced. Printed slips kept literal: no printed word, symbol, hypothesis or number of the counted statement or of its proof was altered. Layout and class: the article's own class is \texttt{article} with packages including geometry, amsmath, amssymb, amsthm, booktabs; the wrapper is the lane's pinned \texttt{amsart} editorial wrapper at 10pt on A4, which restates the theorem-like environments the retained text needs and adds the article's own macro definitions that the retained text needs (\texttt{\textbackslash KS}, \texttt{\textbackslash LB}, \texttt{\textbackslash bg}, \texttt{\textbackslash refi}), with extra wrapper packages (bm, bbm, dsfont, xspace, cancel, mathtools). The delivered numbering is the unit's own. Assets: no retained span contains a figure environment, a graphics-inclusion command, a PDF-page inclusion, a table or a TikZ picture; no image file is copied into this package, the package declares no asset dependency and no image file is redistributed with it. Licence: version arXiv:2607.09555v1 of this article is distributed under the Creative Commons Attribution 4.0 International licence, as recorded in the arXiv licence metadata for this exact version and shown on the abstract page https://arxiv.org/abs/2607.09555v1. A full rights notice is delivered at licenses/arxiv-2607.09555-CC-BY-4.0.txt. Characters: every retained excerpt is pure ASCII, so the delivered engine, XeLaTeX with its default Latin Modern text font, reports no missing character.
\begin{equation}\label{eq:kernel}
\bar v(d) =
\begin{cases}
-1, & d\le 2,\\
V_{\mathrm{LJ}}(d-1), & d\ge 3,
\end{cases}
\end{equation}

\begin{equation}\label{eq:elb}
E_\LB(n) = \sum_{k=1}^D f(n_k) + \sum_{i<j} \bar v(|i-j|)\, n_i n_j,
\end{equation}

\begin{equation}\label{eq:f}
f(n) =
\begin{cases}
V^*_n, & 0\le n\le 6,\\[4pt]
\dfrac{n(n-1)}{30}\, V^*_6, & n\ge 7,
\end{cases}
\qquad
(V^*_{0..6}) = (0,\,0,\,-1,\,-3,\,-6,\,-9.103852,\,-12.712062).
\end{equation}

\begin{equation}\label{eq:elbsplit}
E_\LB = E_\bg(D) - \frac e2 - Q(e) - \sum_i e_i\,\Phi(i),
\end{equation}

\begin{equation}\label{eq:square}
\frac12(e_1^2+e_2^2)+e_1 e_2 = \frac12(e_1+e_2)^2
\end{equation}

\begin{equation}\label{eq:minelb}
\min_n E_\LB = E_\bg(D) - \frac e2 - \frac12 e^2 - e\,\Phi_c.
\end{equation}

\begin{lemma}[Per-layer penalty]\label{lem:penalty}
For every integer $n\ge0$,
\[
\alpha\,(n-4)_+^2 \;\le\; \delta(n) \;\le\; \alpha\,n(n-1),
\]
with equality on the right for $n=6$ and all $n\ge7$. Equivalently, in the surplus variable $x=n-1\ge0$, $\ \delta(1+x)\ge\alpha\,(x-3)_+^2$.
\end{lemma}

\begin{proposition}[Exact asymptotic gap between the two programs]\label{prop:gap}
Fix the kernel $\bar v$ of \eqref{eq:kernel}. For a span of $D\ge3$ layers and surplus $e=N-D\ge0$, let $M_\KS(D,e)$ and $M_\refi(D,e)$ denote the minima of \eqref{eq:elb} over profiles $n_k\ge1$, $\sum_k n_k=N$, with $f(n)=-\binom n2$ and with the $f$ of \eqref{eq:f}, respectively. Then
\[
\frac\alpha3\,e^2 - 6\alpha e
\;\le\;
M_\refi(D,e)-M_\KS(D,e)
\;\le\;
\frac\alpha3\,e^2 + (1+3\alpha)\,e + 3\alpha.
\]
In particular $M_\refi=M_\KS+\frac\alpha3 e^2+O(e)$, uniformly in $D$.
\end{proposition}

\begin{proof}
\emph{Lower bound.} For any admissible profile $n$, with surpluses $e_k=n_k-1\ge0$,
\[
E_\refi(n) = E_\KS(n) + \sum_k \delta(n_k).
\]
By \eqref{eq:elbsplit} and \eqref{eq:minelb},
\[
E_\KS(n)-M_\KS = \Bigl[\tfrac12 e^2 - Q(e)\Bigr] + \Bigl[e\,\Phi_c - \sum_i e_i\Phi(i)\Bigr],
\]
and both brackets are nonnegative. For the first,
\[
\tfrac12 e^2 - Q(e) = \sum_{i<j}\bigl(1-\kappa(|i-j|)\bigr)e_i e_j \;\ge\; \beta \!\!\sum_{|i-j|\ge3}\!\! e_i e_j,
\]
since $\kappa=1$ at distances $1,2$ and $\kappa\le\kappa(3)$ beyond. Partition the layers into the three residue classes $k\bmod3$; distinct layers in one class lie at distance $\ge3$, so with $S_r$, $T_r$ the sum and the sum of squares of the surpluses in class $r$, and $T=\sum_r T_r=\sum_k e_k^2$,
\[
\beta \!\!\sum_{|i-j|\ge3}\!\! e_i e_j \;\ge\; \frac\beta2 \sum_r \bigl(S_r^2 - T_r\bigr).
\]
By Lemma~\ref{lem:penalty} and $(x-3)_+^2\ge x^2-6x$ (valid for all $x\ge0$), $\sum_k\delta(n_k)\ge\alpha\,(T-6e)$. Adding, then using $T\le\sum_r S_r^2$ and $\alpha\le\beta/2$,
\[
E_\refi(n)-M_\KS \;\ge\; \frac\beta2\sum_r S_r^2 + \Bigl(\alpha-\frac\beta2\Bigr)T - 6\alpha e
\;\ge\; \alpha\sum_r S_r^2 - 6\alpha e
\;\ge\; \frac\alpha3\,e^2 - 6\alpha e,
\]
the last step by Cauchy--Schwarz with $\sum_r S_r=e$. Minimizing over profiles gives the lower bound.

\emph{Upper bound.} Let $n^{(3)}$ spread the surplus as evenly as possible over the three most central layers ($D\ge3$). All surplus pairs sit at index distance $\le2$, where $\kappa=1$, so by \eqref{eq:square} $Q(e)=\tfrac12 e^2$ exactly; and since $\Phi(p)-\Phi(p\pm1)=\kappa(D-p)-\kappa(p)$ in absolute value is at most $1$, each of the three sites carries a field within $1$ of $\Phi_c$, so $\sum_i e_i\Phi(i)\ge e(\Phi_c-1)$. Hence, by \eqref{eq:minelb}, $E_\KS(n^{(3)})\le M_\KS+e$. By Lemma~\ref{lem:penalty} and $e_k\le e/3+1$,
\[
\sum_k \delta(n_k) \;\le\; \alpha\sum_k n_k(n_k-1) = \alpha\Bigl(\sum_k e_k^2 + e\Bigr) \;\le\; \alpha\Bigl(\frac{e^2}3+3e+3\Bigr).
\]
Adding the two estimates bounds $M_\refi\le E_\refi(n^{(3)})$ as stated.
\end{proof}

\end{document}
